% !TEX TS-program = xelatex
\documentclass[a4paper,12pt]{article}

% --- Packages principaux ---
\usepackage{fontspec}
\usepackage{xcolor}
\usepackage{graphicx}
\usepackage{geometry}
\usepackage{setspace}
\usepackage{titlesec}
\usepackage{enumitem}
\usepackage{ragged2e}
\usepackage{tikz}

% --- Marges et interlignes ---
\geometry{
  top=2.5cm,
  bottom=2.5cm,
  left=2.5cm,
  right=2.5cm
}
\setstretch{1.1}

% --- Polices modernes ---
\setmainfont{Arial} % tu peux changer pour Helvetica, Times New Roman, etc.

% --- Commandes de style ---
\newcommand{\ligneRouge}{
  \vspace{3mm}
  \noindent\textcolor{red!70!black}{\rule{0.75\textwidth}{0.8pt}}
  \vspace{3mm}
}

% --- En-têtes de niveau personnalisés ---
\titleformat{\section}{\large\bfseries\centering}{}{0pt}{}
\titleformat{\subsection}{\bfseries\centering}{}{0pt}{}

% --- Début du document ---
\begin{document}

% === LOGO UNIVERSITÉ ===
\vspace*{-1cm}
\noindent
%\includegraphics[height=3cm]{logo_universite.png} % %% À compléter : chemin du logo
\hfill
\vspace{1cm}

% === TITRE PRINCIPAL ===
\begin{center}
  {\fontsize{36}{42}\selectfont\textbf{\textcolor{red!70!black}{THÈSE}}}\\[1.2cm]

  {\large En vue de l’obtention du}\\[2mm]
  {\Large\textbf{DOCTORAT DE L’UNIVERSITÉ DE TOULOUSE}}\\[3mm]
  {\normalsize Délivré par l’Université Toulouse 3 – Paul Sabatier}
\end{center}

\vspace{1cm}

% === NOM DU CANDIDAT ===
\begin{center}
  \ligneRouge
  \textbf{Présentée et soutenue par}\\[2mm]
  {\Large\textbf{[Prénom NOM]}}\\[2mm] % %% À compléter
  Le [Date de soutenance]\\[4mm]
  \textbf{Titre de la thèse :}\\[2mm]
  \textit{[Titre complet de la thèse]} % %% À compléter
  \ligneRouge
\end{center}

\vspace{3mm}

% === INFORMATIONS DOCTORALES ===
\begin{center}
  \textbf{École doctorale :} SDM – Sciences de la Matière – Toulouse\\[2mm]
  \textbf{Spécialité :} Physico-Chimie Théorique\\[2mm]
  \textbf{Unité de recherche :}\\
  LCPQ–IRSAMC – Laboratoire de Chimie et Physique Quantiques\\[2mm]
  \textbf{Thèse dirigée par :}\\
  Nicolas SUAÜD et Nathalie GUIHÉRY
\end{center}

\vspace{0.6cm}

% === JURY ===
\begin{center}
  {\large\textbf{Jury}}\\[3mm]
  \begin{tabular}{l l}
    \textbf{Mme [Prénom NOM]} & Rapporteure \\[2mm]
    \textbf{M. [Prénom NOM]} & Rapporteur \\[2mm]
    \textbf{M. [Prénom NOM]} & Examinateur \\[2mm]
    \textbf{Mme [Prénom NOM]} & Examinatrice \\[2mm]
    \textbf{M. [Prénom NOM]} & Directeur de thèse \\[2mm]
    \textbf{Mme [Prénom NOM]} & Co-directrice de thèse \\[2mm]
    \textbf{M. [Prénom NOM]} & Invité \\[2mm]
    \textbf{M. [Prénom NOM]} & Invité \\
  \end{tabular}
\end{center}

\vfill

% === FILIGRANE (optionnel) ===
% Si tu veux ajouter le logo en arrière-plan :
% \begin{tikzpicture}[remember picture,overlay]
%   \node[opacity=0.05,inner sep=0pt] at (current page.center)
%     {\includegraphics[width=0.7\paperwidth]{logo_fond.png}}; % %% À compléter si souhaité
% \end{tikzpicture}

\end{document}
